\part{Hybrid and recurrent phases}

\chapter{Why hybrid token mixers matter}
\label{ch:hybrid-mixers}
Hybrid token mixers explore architectures beyond pure attention.



\begin{intuitionbox}{Intuition}
Attention stores history explicitly; recurrent mixers compress it into state.
\end{intuitionbox}

\begin{definitionbox}{Mathematical or contract statement}
A heterogeneous stack may contain attention blocks, MLPs, and recurrent blocks, each with its own state.
\end{definitionbox}

\begin{implementationbox}{Implementation interpretation}
Let layers declare state kind: KV cache, recurrent state, both, or neither.
\end{implementationbox}

\begin{verificationbox}{How to verify this works}
\begin{itemize}
\item Check shapes and dtypes at the boundary.
\item Compare deterministic outputs against PyTorch golden vectors.
\item Record tolerance, seed, model artifact, and backend.
\item Do not benchmark until validation passes.
\end{itemize}
\end{verificationbox}

\begin{warningbox}{Common failure modes}
\begin{itemize}
\item Letting framework defaults choose dtype or layout.
\item Testing only a batch size of one.
\item Depending on upstream checkpoint names in runtime code.
\item Treating benchmark output as validation.
\end{itemize}
\end{warningbox}

\begin{chapterexercises}
\item State the central invariant in one sentence.\\
\textit{Hint.} Use the conceptual guide vocabulary.
\item Construct a toy example with small shapes.\\
\textit{Hint.} Use $B=2$, $T=3$, $H=8$ where possible.
\item Name a validation test that would catch a plausible bug.\\
\textit{Hint.} Prefer generated golden vectors to handwritten long arrays.
\end{chapterexercises}

\chapter{State-space and recurrent intuition}
\label{ch:state-space}
State-space and recurrent models update hidden state as tokens arrive. \citep{gu2024mamba}

\[
h_t=A_t h_{t-1}+B_t x_t,\qquad y_t=C_t h_t+D_t x_t,
\]
where \(h_t\) is recurrent state, \(x_t\) is the token representation, and the matrices may be input-dependent in selective state-space models.

\begin{intuitionbox}{Intuition}
A recurrent state is a summary with an update law.
\end{intuitionbox}

\begin{definitionbox}{Mathematical or contract statement}
At a high level, a recurrent token mixer maps \((h_{t-1},x_t)\) to \((h_t,y_t)\). Sequence-mode evaluation and repeated step-mode evaluation must produce matching \(y_{1:T}\) within tolerance.
\end{definitionbox}

\begin{implementationbox}{Implementation interpretation}
Validate sequence scan against repeated step updates.
\end{implementationbox}

\begin{verificationbox}{How to verify this works}
\begin{itemize}
\item Check shapes and dtypes at the boundary.
\item Compare deterministic outputs against PyTorch golden vectors.
\item Record tolerance, seed, model artifact, and backend.
\item Do not benchmark until validation passes.
\end{itemize}
\end{verificationbox}

\begin{warningbox}{Common failure modes}
\begin{itemize}
\item Letting framework defaults choose dtype or layout.
\item Testing only a batch size of one.
\item Depending on upstream checkpoint names in runtime code.
\item Treating benchmark output as validation.
\end{itemize}
\end{warningbox}

\begin{chapterexercises}
\item State the central invariant in one sentence.\\
\textit{Hint.} Use the conceptual guide vocabulary.
\item Construct a toy example with small shapes.\\
\textit{Hint.} Use $B=2$, $T=3$, $H=8$ where possible.
\item Name a validation test that would catch a plausible bug.\\
\textit{Hint.} Prefer generated golden vectors to handwritten long arrays.
\end{chapterexercises}

\chapter{Mamba-style blocks}
\label{ch:mamba-blocks}
Mamba-style blocks are a concrete later-phase recurrent or state-space family. \citep{gu2024mamba}



\begin{intuitionbox}{Intuition}
A Mamba-style block replaces all-pairs attention with a structured scan.
\end{intuitionbox}

\begin{definitionbox}{Mathematical or contract statement}
The exact parameterisation is model-specific, but the engine abstraction can be common.
\end{definitionbox}

\begin{implementationbox}{Implementation interpretation}
Add \code{MambaBlock} with state allocation, scan, step update, and exports.
\end{implementationbox}

\begin{verificationbox}{How to verify this works}
\begin{itemize}
\item Check shapes and dtypes at the boundary.
\item Compare deterministic outputs against PyTorch golden vectors.
\item Record tolerance, seed, model artifact, and backend.
\item Do not benchmark until validation passes.
\end{itemize}
\end{verificationbox}

\begin{warningbox}{Common failure modes}
\begin{itemize}
\item Letting framework defaults choose dtype or layout.
\item Testing only a batch size of one.
\item Depending on upstream checkpoint names in runtime code.
\item Treating benchmark output as validation.
\end{itemize}
\end{warningbox}

\begin{chapterexercises}
\item State the central invariant in one sentence.\\
\textit{Hint.} Use the conceptual guide vocabulary.
\item Construct a toy example with small shapes.\\
\textit{Hint.} Use $B=2$, $T=3$, $H=8$ where possible.
\item Name a validation test that would catch a plausible bug.\\
\textit{Hint.} Prefer generated golden vectors to handwritten long arrays.
\end{chapterexercises}

\chapter{Recurrent state versus KV-cache}
\label{ch:state-vs-cache}
KV cache and recurrent state are both generation state but have different semantics.



\begin{intuitionbox}{Intuition}
The type of state determines the equivalence law.
\end{intuitionbox}

\begin{definitionbox}{Mathematical or contract statement}
KV cache equivalence compares full attention recomputation with cached decode; recurrent equivalence compares scan with repeated step updates.
\end{definitionbox}

\begin{implementationbox}{Implementation interpretation}
Use separate tagged state variants.
\end{implementationbox}

\begin{verificationbox}{How to verify this works}
\begin{itemize}
\item Check shapes and dtypes at the boundary.
\item Compare deterministic outputs against PyTorch golden vectors.
\item Record tolerance, seed, model artifact, and backend.
\item Do not benchmark until validation passes.
\end{itemize}
\end{verificationbox}

\begin{warningbox}{Common failure modes}
\begin{itemize}
\item Letting framework defaults choose dtype or layout.
\item Testing only a batch size of one.
\item Depending on upstream checkpoint names in runtime code.
\item Treating benchmark output as validation.
\end{itemize}
\end{warningbox}

\begin{chapterexercises}
\item State the central invariant in one sentence.\\
\textit{Hint.} Use the conceptual guide vocabulary.
\item Construct a toy example with small shapes.\\
\textit{Hint.} Use $B=2$, $T=3$, $H=8$ where possible.
\item Name a validation test that would catch a plausible bug.\\
\textit{Hint.} Prefer generated golden vectors to handwritten long arrays.
\end{chapterexercises}

\chapter{Interface design for heterogeneous layers}
\label{ch:heterogeneous-layers}
Heterogeneous layers require an interface that declares parameters, state kind, prefill behaviour, decode behaviour, and validation exports.



\begin{intuitionbox}{Intuition}
Heterogeneity should be explicit, not hidden in a switch statement.
\end{intuitionbox}

\begin{definitionbox}{Mathematical or contract statement}
A layer maps \((x,state)\) to \((y,state')\) and may expose diagnostics.
\end{definitionbox}

\begin{implementationbox}{Implementation interpretation}
Keep forward, state allocation, state clear, and profiling labels narrow.
\end{implementationbox}

\begin{verificationbox}{How to verify this works}
\begin{itemize}
\item Check shapes and dtypes at the boundary.
\item Compare deterministic outputs against PyTorch golden vectors.
\item Record tolerance, seed, model artifact, and backend.
\item Do not benchmark until validation passes.
\end{itemize}
\end{verificationbox}

\begin{warningbox}{Common failure modes}
\begin{itemize}
\item Letting framework defaults choose dtype or layout.
\item Testing only a batch size of one.
\item Depending on upstream checkpoint names in runtime code.
\item Treating benchmark output as validation.
\end{itemize}
\end{warningbox}

\begin{chapterexercises}
\item State the central invariant in one sentence.\\
\textit{Hint.} Use the conceptual guide vocabulary.
\item Construct a toy example with small shapes.\\
\textit{Hint.} Use $B=2$, $T=3$, $H=8$ where possible.
\item Name a validation test that would catch a plausible bug.\\
\textit{Hint.} Prefer generated golden vectors to handwritten long arrays.
\end{chapterexercises}

\chapter{Validation for hybrid models}
\label{ch:hybrid-validation}
Hybrid validation extends transformer validation with state-update checks.



\begin{intuitionbox}{Intuition}
Each new state semantics needs its own equivalence law.
\end{intuitionbox}

\begin{definitionbox}{Mathematical or contract statement}
Validate component outputs, full logits, recurrent scan equivalence, attention cache equivalence, and greedy ids.
\end{definitionbox}

\begin{implementationbox}{Implementation interpretation}
Start with tiny synthetic configs and generated vectors.
\end{implementationbox}

\begin{verificationbox}{How to verify this works}
\begin{itemize}
\item Check shapes and dtypes at the boundary.
\item Compare deterministic outputs against PyTorch golden vectors.
\item Record tolerance, seed, model artifact, and backend.
\item Do not benchmark until validation passes.
\end{itemize}
\end{verificationbox}

\begin{warningbox}{Common failure modes}
\begin{itemize}
\item Letting framework defaults choose dtype or layout.
\item Testing only a batch size of one.
\item Depending on upstream checkpoint names in runtime code.
\item Treating benchmark output as validation.
\end{itemize}
\end{warningbox}

\begin{chapterexercises}
\item State the central invariant in one sentence.\\
\textit{Hint.} Use the conceptual guide vocabulary.
\item Construct a toy example with small shapes.\\
\textit{Hint.} Use $B=2$, $T=3$, $H=8$ where possible.
\item Name a validation test that would catch a plausible bug.\\
\textit{Hint.} Prefer generated golden vectors to handwritten long arrays.
\end{chapterexercises}
